\chapter*{Exercices de révision}

\section*{Calcul algébrique}


\begin{itemize}
\it On donne sur un axe quatre points $A$, $B$, $C$, $D$. Démontrer la relation : $\overline{AC} + \overline{BD} = \overline{AD} + \overline{BC}$.
\it Soient sur un axe deux points $A$ et $B$ d'abscisses $a$ et $b$. Calculer les abscisses des points $M$ et $N$ qui partagent $[AB]$ en $3$ parties égales.
\it On donne sur un axe trois points $A$, $B$ et $C$ d'abscisses $a$, $b$ et $c$. Trouver l'abscisse du point $G$ tel que $\overline{GA}+\overline{GB}+\overline{GC}=0$.
\it Soient sur un axe $4$ points, $A$, $B$, $C$ et $D$. Démontrer la relation : \[ \overline{AB}.\overline{CD} + \overline{BC}.\overline{AD}+\overline{CA}.\overline{BD}=0.\]
\it On donne sur un axe $3$ points, $A$, $B$ et $M$. On désigne par $O$ le milieu de $[AB]$. Démontrer les relations : 
\[ \overline{MA} + \overline{MB} = 2 \overline{MO}.\phantom{miaoumiaoumiaou} \overline{MA}.\overline{MB}=
\overline{MO}^2-\frac{\overline{AB}^2}4.\]
\it Soient deux points $A$ et $B$ d'un axe, $O$ le milieu de $[AB]$ et $M$ un point quelconque de l'axe.
Démontrer les relations : 
\[ \overline{MA}^2 + \overline{MB}^2 = 2 \overline{MO}^2+\frac{\overline{AB}^2}2.\phantom{miaoumiaoumiaou} \overline{MA}^2-\overline{MB}^2=2\overline{AB}.\overline{OM}\]
\item Effectuer les produits suivants : \begin{itemize}
\it $\left(\frac34 a^2x^3y^2\right)\left(-\frac29 bx^2y^4\right)(-10a^2b^3x^4y^2)$.
\it $(2x+3)(x-5)(3x+1)$.
\it $(2x^3-3x^2+4x)(4x^2+6x+8)$.
\it $(2x-3y)^2(x-y+1)$.
\end{itemize}
\item Calculer les expressions suivantes : \begin{itemize}
\it $x(x+1)(x+2)-3(x-3)(x+2)+2(x-6)$.
\it $(2x+3y)(2y-1)-(2x-y)(5x-1)+(2x-3y)(5x+2y)$.
\it $(5x+3y)(3x-2y)+(5x+2y)(3y+1)-(5x-2y)(3x+1)$.
\it $(ax-y)(ax-2y)+3y(ax-y+1)-y(8y+3)$.
\end{itemize}
\item Établir les identités suivantes : 
\begin{itemize}
\it $(a+b-c)^2+(a-b+c)^2=2a^2+2(b-c)^2$.
\it $(a+b+c)^2-(a-b-c)^2=4a(b+c)$.
\it $(a+b)^3+(a-b)^3 = 2a(a^2+3b^2)$.
\it $(ab+1)^2+(a-b)^2=(a^2+1)(b^2+1)$.
\it $(ab+1)^2-(a+b)^2=(a^2-1)(b^2-1)$.
\end{itemize}
\item Décomposer en un produit de facteurs les expressions suivantes : 
\begin{itemize}
\it $12a^2b-3b^3$.
\it $48x^2y^2-27y^4$.
\it $x^2+6x+9$. 
\it $3x^2-30x+75$.
\it $(2x+3)^2-49$.
\it $x^2-10x+21$.
\it $(2a-b)^2-a^2$.
\it $(2a+b)^2-(a+2b)^2$.
\it $(ax+1)^2-(x+a)^2$. 
\it $(ax+b)^2+(bx-a)^2$.
\end{itemize}
\item Simplifier les expressions suivantes : 
\begin{itemize}
\it $\frac{(a+b)^2-c^2}{(a+c)^2-b^2}$.
\it $\frac{(3a+2b)^2-(a+2b)^2}{a^2-b^2}$.
\it $\frac{x^2-2x+1}{x^2-1}$.
\it $\frac{4x^2-9}{4x^2+8x-21}$.
\it $\frac{a+1}{a-1}+\frac{a-1}{a+1}-\frac{4a}{a^2-1}$.
\it $\frac2{2a+b}-\frac1{2a-b}+\frac{2b}{4a^2-b^2}$.
\it $\frac{x^2}{x-1}+\frac1{x+1}-\frac{2x}{x^2-1}$.
\it $\frac{x+1}x-\frac2{x^2+2x}+\frac{x-1}{x+2}$.
\it $\frac{\frac{x+1}x-\frac{x}{x+1}}{1+\frac{x}{1+x}}$.
\it $\frac{\frac{a-b}b+\frac{2b}{a+b}}{\frac{a+b}a-\frac{2b}{a+b}}$.
\end{itemize}
\it Soit les rapports égaux $\frac{a}{a'}=\frac{b}{b'}=\frac{c}{c'}$. Démontrer que chacun de ces rapports est égal à $\frac{5a-3b+2c}{5a'-3b'+2c'}$.
\it Soit l'égalité $\frac{a}b=\frac{c}d$. Démontrer que $\frac{a^2+b^2}{c^2+d^2} = \frac{ab}{cd}$.
\it Démontrer que si $a$, $b$, $c$, $d$ vérifient la relation $(ad+bc)^2=4abcd$, ces quatre nombres vérifient l'égalité $ad=bc$. 
\it Démontrer que si $(a^2+b^2)(c^2+d^2) = (ac+bd)^2$, les quatre nombres $a$, $b$, $c$ et $d$ vérifient l'égalité $ad=bc$.
\end{itemize}

\section*{Équations du premier degré}
\begin{itemize}
\item Résoudre les équations suivantes : 
\begin{itemize}
\it $\frac{x-1}3+\frac{4(2x-3)}9 = \frac{5x+1}8$.
\it $\frac{x+5}6 + \frac{x+9}8 = \frac{2x+7}9$.
\it $\frac{5x-9}3 + \frac{7x-5}9 = \frac{3x-143}4$.
\it $\frac{13x-16}{15}+\frac{x-32}{35} = \frac{295-x}21$.
\it $\frac{3x-2}{18}+\frac{2x-9}{45}=\frac{x+3}{10}$.
\it $\frac{4x+9}{77}-\frac{5x+12}{14}=\frac{11x+28}{22}$.
\it $\frac{x+5}4 -\frac{x-3}6 = \frac{x}3$.
\it $-\frac{x+1}3 = \frac{3x+4}5$.
\it $\frac{5x-7}4 + \frac{3x-5}8 = \frac{9x-4}5$.
\it $\frac{4x-1}3+\frac{2x-1}2=\frac{10x+63}7$.
\it $\frac{x-8}5-\frac{x-23}{15} = \frac{43-x}{18} $.
\it $\frac{13x-34}{21} - \frac{3x-7}{56} = \frac{19-7x}{24}$. 
\end{itemize}
\item Résoudre les équations suivantes : 
\begin{itemize}
\it $\frac{3x-1}7 + \frac{3x-4}{5x-3} = \frac{6x+5}{14}$. 
\it $\frac{2x-7}{2(x-1)} = 1+ \frac1{4-x}$.
\it $\frac3{x-3} + \frac4{x-5} = \frac7{x-4}$.
\it $\frac{2x+3}{5x-3} = \frac{3}{2(2x-3)} + \frac25$.
\it $\frac{9x+2}{6(x+1)} - \frac{5x+8}{2x+1} = - \frac{x+5}{x+1}$. 
\it $\frac5{2(3x+2)} - \frac1{2x+3} = \frac1{3x-8}$. 
\end{itemize}
\item Résoudre les équations suivantes : 
\begin{itemize}
\it $(2x-3)(4x-12)=0$.
\it $x(x-7)(3x-2) = 0$.
\it $x^2-(4x+5)^2 = 0$. 
\it $(2x-3)^2- (x+6)^2=0$.
\it $\frac4{3-2x}+\frac9{x+3} = \frac{13}{2x+3}$.
\it $\frac4{x+10}+\frac5{x-20} = \frac9{x+20}$.  
\end{itemize}
\item Résoudre les équations suivantes où le nombre $m$ est supposé connu : 
\begin{itemize}
\it $\frac{8x-5m}7 - \frac{5x-11m}5 = \frac{78 - x}{35}$. 
\it $m(x-m)=x-1$.
\it $\frac{4x+m}5 - \frac34(x-1) = \frac{2x-m+3}{10}$.
\it $x+m^2+1=m(x+2)$.
\end{itemize}
\it Soit l'équation : \[ \frac{5x-m}3 - \frac{4x+m}5
= \frac{7(3x+4)}{14}\]
dans laquelle $m$ est un nombre connu.
\begin{enumerate}
\item Résoudre cette équation.
\item Déterminer $m$ de façon que $x=1$.
\item Calculer les valurs de $m$ pour lesquelles $x>3$.
\end{enumerate}
\end{itemize}
\section*{Problèmes du premier degré}
\begin{itemize}
\it La somme de deux nombres est $324$. En ajoutant $26$ à chacun d'eux, l'un devient le triple de l'autre. Quels sont ces deux nombres ? 
\it La différence de deux nombres est $1~032$. Le quotient entier de ces deux nombres est $13$ et le reste de leur division est $48$. Quels sont ces deux nombres ? 
\it Deux cyclistes partent en même temps de deux villes A et B distantes de $200$ km et vont à la rencontre l'un de l'autre. Ils se rencontrent au bout de $4$ heures. Si le cycliste qui part de A était parti une demi-heure avant l'autre, la rencontre aurait eu lieu $3$ heure $48$ minutes après le départ du deuxième cycliste. Quelle est la vitesse de chacun d'eux ? 
\it Un train a mis $36$ secondes à passer devant un observateur immobile. Sa longueur est $300$ m. Quelle est sa vitesse ? Un second train met $24$ secondes pour croiser le premier et passe devant l'observateur immobile en $18$ secondes. Trouver la longueur et la vitesse du second train. 
\it Deux cyclistes sont séparés par une distance de $54$ km. S'ils allaient à la rencontre l'un de l'autre, ils se rencontreraient au bout d'une heure. S'ils allaient dans le même sens, ils se rejoindraient au bout de $9$ heures. Quelle est la vitesse de chaque cycliste ? 
\it Deux capitaux ont pour somme $3~000$ F.L'un est placé à 6\% et l'autre à 5\%. La somme de leurs intérêts annuels est 162 F. Quelle est la valeur de chaque capital ? 
\it Deux lingots d'or sont l'un au titre de 0,95, l'autre au titre de 0,8. On les fond ensemble en ajoutant 3 kg d'or pur. Le lingot ainsi obtenu est au titre de 0,906 et pèse 37,5 kg. Quel est le poids respectif des lingots primitifs. 
\it On veut obtenir 450 g d'un alliage d'or au titre de 0,83 en fondant des lingots aux titres respectifs de 0,8 et 0,9. Quels doivent être les poids de ces deux lingots ? 
\it Un épicier achète 100 kg de café, une partie à 9,60 F le kilogramme, l'autre à 11,20 F le kilogramme. Il revend le tout pour 1~209,60 F avec un bénéfice de 20\% sur le prix d'achat. Quels sont les poids respectifs de chaque qualité de café ? 
\it Un terrain rectangulaire a 414 m de périmètre. Si la longueur augmentait du tiers de sa valeur et si la largeur diminuait du tiers de la sienne, le périmètre augmenterait de 26 m. Quelles sont les dimensions de ce  terrain ? 
\it La longueur d'un rectangle est supérieure à sa largeur de 21 m. Si on augmentait ses dimensions de 4 m, la surface augmenterait de 492 $m^2$.Trouver les dimensions de ce rectangle. 
\it Deux capitaux sont tels que le second surpasse de 2~300 F les $\frac34$ du premier. En les plaçant à 3,6\%, le premier pendant 7 mois, le second pendant 5 mois, l'intérêt du premier dépasse de 207,30 F l'intérêt du second. Calculer les deux capitaux. 
\it Deux villes A et B sont distantes de 600 km. La tonne de charbon coûte 120 F en A et 144 F en B; et le transport coûte 0,15 F par tonne et par kilomètre. Trouver le point situé entre A et B où le charbon revient au même prix, soit qu'il vienne de A, soit qu'il vienne de B. 
\it Une somme de 2~170 F est partagée entre deux personnes. La première ayant dépensé les $\frac57$ de sa part et la deuxième les $\frac25$ de la sienne, il leur reste la même sommme. Trouver ces deux parts. 
\it Deux cyclistes partent en même temps de deux villes A et B et vont à la rencontre l'un de l'autre. Celui qui part de A fait 6 km à l'heure de plus que l'autre et s'arrête un quart d'heure après chaque heure de marche. L'autre, qui part de B, n'a qu'un seul arrêt de 12 minutes. Les cyclistes se rencontrent au bout de 5 heures au milieu de [AB]. Trouver la distance AB et la vitesse de chaque cycliste. 
\it Partager une somme de 41~340 F proportionnellement aux trois nombres 11, 13 et 15. 
\it Les fortunes de 3 personnes sont proportionnelles aux nombres 2, 3 et 4. En additionnant la première, le double de la seconde et le triple de la troisième, on trouve 1~000 F. Quels sont ces trois fortunes ? 
\it Un cultivateur a vendu 2~880 F sa récolte de blé et 1~440 F sa récolte d'avoine. Le poids du blé surpasse de 16 quintaux celui de l'avoine. Sachant que le prix du quintal de blé est les $\frac85$ de celui du quintal d'avoine : 
\begin{enumerate}
\item Trouver les quantités vendues de chacune des deux cérales. 
\item Quel est le prix du quintal de chacune d'elles ?
\end{enumerate}
\it Un automobiliste part à 8h40 d'une ville A et se rend à une ville B située à 39 km de A. Au retour, il fait un détour de 21 km, mais sa vitesse horaire surpasse de 10 km à l'heure sa vitesse à l'aller, si bien que la durée du trajet retour est égale aux $\frac43$ de la durée du trajet aller.\begin{enumerate}
\item Calculer la vitesse de l'automobiliste à l'aller et au retour.
\item Il s'est arrêté 1h 26 min en B. Quelle est l'heure de son retour en A ? 
\end{enumerate}
\it Un chemisier a acheté deux lots de chemises identiques. Il vend le premier lot pour 576 F en réalisant un bénéfice de 2,40 F l'unité. Il vend une partie du second lot pour 180 F avec un bénéfice de 3F et le reste pour 252 F avec un bénéfice de 2 F par chemise. Sahcant que le nombre de chemises du deuxième lot est les $\frac34$ de celui du premier lot : \begin{enumerate}
\item Trouver le prix d'achat d'une chemise.
\item Trouver le nombre de chemises vendues à chaque fois.
\end{enumerate}
\it Un cycliste effectue un trajet comprenant 36 km de terrain plat, 24 km de montées et 48 km de descentes. Dans les montées, sa vitesse horaire moyenne diminue de 12 km et en descente, elle augmente de 15 km. Sachant que al durée du trajet en terrain plat est le tiers de la durée totale du trajet : \begin{enumerate}
\item Trouver la vitesse horaire du cycliste en terrain plat.
\item Quelle est la durée totale du trajet et la vitesse moyenne réalisée ?
\end{enumerate}
\it Un automobiliste part à 8h20 d'une localité A pour une ville B située à 192 km de A et où il doit arriver à une heure déterminée. Il calcule la vitesse horaire moyenne qu'il doit réaliser, mais arrivé à mi-route, il s'aperçoit que sa vitesse horaire a été inférieur de 12 km à la vitesse prévue. Il accélère son allure et arrive à l'heure fixée, en effectuant dans la deuxième partie de son parcours une vitesse horaire supérieure de 18 km à la vitesse prévue. \begin{enumerate}
\item Calculer la vitesse moyenne réalisée sur le parcours entier, puis l'heure d'arrivée. 
\item Quel était le retard à mi-route de l'automobiliste ?
\end{enumerate}
\end{itemize}